\section*{Preface}
\vspace{-2ex}\titlerule\vspace{4ex}

This thesis was carried out as a degree project for the Master of Engineering Science with a specialization in AI and Automation at the Department of Engineering Science, University West, Sweden. The work was conducted in collaboration with AP\&T, a company specializing in automated production lines for the automotive industry, during the spring semester of 2026.

The project addresses the automation of pyrometer data pre-processing for metal forming and heat-treatment processes. This individual thesis contribution focuses specifically on sensor calibration~(ATP-2) and additional data compression methods~(ATP-3), investigating both classical regression methods and machine learning-based approaches for the calibration stage, and Delta Encoding, Variational Autoencoder, and Deep Autoencoder approaches for the compression stage. The work shares a common base pipeline and dataset with a parallel thesis contribution by Mallepalli Sravya Reddy, which addresses signal denoising~(ATP-1) and classical and PCA-based compression methods. Both contributions are clearly distinguished by their individual research questions, methodology, implementation, and evaluation, in accordance with the requirements confirmed by the examiner.

The NIST open-source additive manufacturing dataset~(Layers~01--10) was used as the primary development and evaluation testbed.

I would like to express sincere gratitude to my supervisors, Amit Kumar Mishra for their continuous guidance, feedback, and support throughout the entire project. I would also like to thank my examiner, Fredrik Sikstr\"{o}m, for the constructive feedback provided during the project.

Special thanks to Karthikeyan Thalavai Pandian~(AP\&T) for providing the two-colour pyrometer calibration formula and for confirming the physical validity of the measured melt pool temperatures.

Finally, I would like to thank my family and friends for their patience and encouragement throughout my studies.

\vspace{1cm}
Trollh\"{a}ttan, Sweden, April 2026

\vspace{1cm}
Avula Ajay Kumar
